%% Lista de acrônimos — manual (a classe usa acronym[printonlyused],
%% que exigiria \ac{} no corpo; o texto expande no primeiro uso, e esta
%% lista serve de referência rápida, no padrão das dissertações do CIn).
\chapter*{Lista de Acrônimos}
\thispagestyle{empty}
\begin{description}\itemsep2pt
  \item[AUC] \emph{Area Under the ROC Curve} — área sob a curva ROC
  \item[CD] \emph{Critical Difference} — distância crítica (Nemenyi)
  \item[CRPS] \emph{Continuous Ranked Probability Score}
  \item[CV] \emph{Cross-Validation} — validação cruzada
  \item[DL] \emph{Deep Learning} — aprendizado profundo
  \item[FM/TFM] \emph{(Tabular) Foundation Model}
  \item[GBDT] \emph{Gradient-Boosted Decision Trees}
  \item[HPO] \emph{Hyperparameter Optimization}
  \item[ICL] \emph{In-Context Learning} — aprendizado em contexto
  \item[KAN] \emph{Kolmogorov--Arnold Network}
  \item[LODO] \emph{Leave-One-Dataset-Out}
  \item[LWTA] \emph{Local Winner-Takes-All}
  \item[MAE] \emph{Mean Absolute Error} — erro absoluto médio
  \item[MLP] \emph{Multi-Layer Perceptron}
  \item[OOF] \emph{Out-Of-Fold} — fora da dobra
  \item[PFN] \emph{Prior-Data Fitted Network}
  \item[PICP] \emph{Prediction Interval Coverage Probability}
  \item[QP] Pergunta de Pesquisa
  \item[QRF] \emph{Quantile Regression Forest} — floresta quantílica
  \item[RMSE] \emph{Root Mean Squared Error} — raiz do erro quadrático médio
  \item[ROPE] \emph{Region of Practical Equivalence} — região de equivalência prática
  \item[SCM] \emph{Structural Causal Model} — modelo causal estrutural
\end{description}
\cleardoublepage
